\section{Design and Implementation}

\begin{table*}[t]
\centering
\small
\caption{
Comparison of OPD supervision interfaces and framework properties.
\cmark\ denotes native or first-class support, \pmark\ denotes partial
or experimental support, and \xmark\ denotes no documented support.
\emph{On-policy self-distillation} refers to settings in which the
teacher is derived from the student model and supervises
student-generated rollouts. The comparison reflects publicly documented
functionality as of July 2026, using the versions or commits
listed in Appendix~\ref{app:framework-comparison}, where per-cell evidence
is also provided.
}
\setlength{\tabcolsep}{6pt}
\begin{tabular}{l cccc ccc}
\toprule
 & \multicolumn{4}{c}{\textbf{OPD supervision-interface coverage}}
 & \multicolumn{3}{c}{\textbf{Framework properties}} \\
\cmidrule(lr){2-5}\cmidrule(lr){6-8}

\textbf{Framework}
 & \shortstack{Logit\\OPD}
 & \shortstack{Cross-tok.\\OPD}
 & \shortstack{On-policy\\self-distill.}
 & \shortstack{Step-wise\\OPD}
 & \shortstack{Released\\YAML/CLI}
 & \shortstack{Documented method\\extension API}
 & \shortstack{Compatible w/\\multi-node} \\
\midrule

TRL
 & \cmark
 & \pmark
 & \pmark
 & \xmark
 & \cmark
 & \xmark
 & \pmark \\

verl
 & \cmark
 & \xmark
 & \xmark
 & \pmark
 & \cmark
 & \xmark
 & \cmark \\

slime
 & \cmark
 & \xmark
 & \pmark
 & \pmark
 & \cmark
 & \xmark
 & \cmark \\

KDFlow
 & \cmark
 & \cmark
 & \cmark
 & \xmark
 & \cmark
 & \pmark
 & \pmark \\

\textbf{Ours}
 & \cmark
 & \cmark
 & \cmark
 & \cmark
 & \cmark
 & \cmark
 & \cmark\,{\footnotesize(via verl)} \\

\bottomrule
\end{tabular}

\label{tab:opd-frameworks}
\end{table*}

% Figure 2: EasyOPD running example (SimCT). All field names, the launch
% command, and the diagnostic metric names are copied from the released
% repository (easyopd/config/simct.yaml, examples/simct/run_simct.sh, and
% easyopd/methods/simct/losses.py).
\begin{figure*}[t]
    \centering
    \footnotesize
    \setlength{\fboxsep}{4pt}
    %---------------- Left: YAML + CLI ----------------%
    \begin{minipage}[t]{0.31\textwidth}
        \raggedright
        \textbf{(a) Declarative configuration}\\[2pt]
        \begin{lstlisting}[basicstyle=\ttfamily\tiny,frame=single,numbers=none,language=]
# config/simct.yaml
student_model_path: Qwen/Qwen3-8B
teacher_model_path:
  meta-llama/Llama-3.2-3B-Instruct
distillation:
  enabled: true
  distillation_loss:
    loss_mode: simct
    cross_tokenizer_kl_direction:
      reverse
    use_task_rewards: false
        \end{lstlisting}
        \textbf{Launch}\\[2pt]
        \begin{lstlisting}[basicstyle=\ttfamily\tiny,frame=single,numbers=none,language=bash]
bash examples/simct/run_simct.sh
        \end{lstlisting}
    \end{minipage}
    \hfill
    %---------------- Middle: dispatch ----------------%
    \begin{minipage}[t]{0.32\textwidth}
        \raggedright
        \textbf{(b) Resolution and hook composition}\\[6pt]
        \fbox{\parbox{0.92\linewidth}{\centering Config}}\\[2pt]
        \centering$\downarrow$\\[2pt]
        \fbox{\parbox{0.92\linewidth}{\centering Registry (\texttt{loss\_mode=simct})}}\\[2pt]
        $\downarrow$\\[2pt]
        \fbox{\parbox{0.92\linewidth}{\centering SimCT module:\\ cross-tokenizer alignment\\ +\ teacher sidecar\ +\ loss}}\\[2pt]
        $\downarrow$\\[2pt]
        \fbox{\parbox{0.92\linewidth}{\centering verl: rollout + optimization}}\\[6pt]
        \raggedright Only the hooks required by SimCT are activated; other methods activate a different subset.
    \end{minipage}
    \hfill
    %---------------- Right: diagnostics ----------------%
    \begin{minipage}[t]{0.31\textwidth}
        \raggedright
        \textbf{(c) Supervision diagnostics}\\[2pt]
        \begin{lstlisting}[basicstyle=\ttfamily\tiny,frame=single,numbers=none,language=]
distillation/overlap_vocab_size
distillation/simct_valid_segments
distillation/simct_span_segments
distillation/simct_skipped_samples
distillation/simct_valid_span_ratio
distillation/simct_loss
        \end{lstlisting}
        \raggedright In addition to task performance, EasyOPD logs whether the intended cross-tokenizer supervision is active.
    \end{minipage}
    \caption{A running example of launching SimCT with \textsc{EasyOPD}. The YAML configuration selects the method and resolves its method-specific options; the registry activates the required alignment, teacher-side, and loss hooks; verl executes rollout generation and optimization, while \textsc{EasyOPD} reports supervision-specific diagnostics. Field names, the launch command, and the diagnostic metric names are taken from the released repository.}
    \label{fig:running_example}
\end{figure*}


EasyOPD provides a unified OPD training process on top of verl. Instead of treating each OPD method as an isolated training script, it decomposes an OPD run into four reusable stages: method invocation, configuration resolution, supervision construction, and distributed execution. This section describes how these stages form a unified workflow that connects diverse OPD methods to the same verl execution substrate.

\subsection{A Unified OPD Workflow}

Although OPD methods differ in their supervision interfaces, they share a common training flow. Given a teacher model, a student model, and training data, the student first generates on-policy rollouts; the selected method then constructs supervision on these rollouts, such as teacher logits, aligned token distributions, judge rewards, or step-level feedback; finally, the training system consumes this supervision through loss computation, reward aggregation, or trajectory-level weighting. EasyOPD makes this shared process explicit and standardizes invocation and dispatch, without requiring all methods to share a single loss, teacher interface, or evaluation metric. As shown in Figure~\ref{fig:framework}, the user layer specifies what to run, the EasyOPD layer defines how supervision is constructed, and the verl layer executes the resulting training process on top of its distributed runtime~\citep{sheng2024hybridflow_verl}.

\subsection{Declarative Invocation and Configuration}

EasyOPD exposes OPD methods through declarative configurations and launch scripts. Released methods can be selected through method-specific YAML configurations without manually editing rollout workers, reward managers, logit processors, or trainer-side loss assembly, which turns method selection into a configuration-level operation and makes experiments easier to reproduce. Figure~\ref{fig:running_example} shows a concrete running example in which the same \texttt{from\_hparams} entry point and a single launch command drive all three studied settings---cross-tokenizer OPD, on-policy self-distillation, and step-wise OPD---together with their baselines, changing only the method name and its YAML configuration.

Different OPD methods require different options: logit-based methods may specify divergence directions or top-$k$ support; cross-tokenizer methods may specify alignment or overlap-vocabulary behavior; rubric-based methods may specify judge prompts or reward aggregation rules; step-wise methods may specify trajectory segmentation or weighting strategies. EasyOPD keeps these choices in a unified configuration space and resolves them through method registration, so new methods can be added without hard-coding method choices into the trainer.

\subsection{Method-Local Supervision Logic}

The core of EasyOPD is to localize OPD-specific supervision logic: a method defines how supervision is obtained, transformed, consumed, and diagnosed. Standard OPD may only require a token-level distillation objective, cross-tokenizer OPD may require tokenizer or span alignment and a common supervision space, rubric-based OPD may convert judge feedback into rewards, and step-wise OPD may attach feedback to intermediate reasoning states or multi-turn trajectories. EasyOPD does not force these methods into a single loss-only abstraction; it gives each a local space to implement the supervision interface it needs.

Method-local logic also includes diagnostics. Implemented method modules expose supervision-specific diagnostics in addition to task performance, so users can check whether the intended supervision signal is active, not only the final loss. For example, the released cross-tokenizer implementation reports the overlap-vocabulary size, the number of valid and span-level aligned segments, the count of skipped samples, and the valid-span ratio, while the step-wise implementation reports the mean step-level weight and the number of supervised steps. These metrics help users debug failures that final task performance alone cannot explain.

\subsection{Execution through Thin verl Hooks}

EasyOPD reuses verl as the execution substrate instead of reimplementing distributed training. verl provides the system-heavy components required by OPD, including rollout generation, reward computation, model forwarding, logit processing, distributed workers, trainer-side optimization, and parallel execution, and EasyOPD connects method-local supervision logic to this substrate through a small set of thin hooks.

A hook is a dispatch boundary rather than a method implementation: it exposes a stable interaction point in the training pipeline and routes the corresponding data or control flow to the selected method module. The hook set covers the main entry points of OPD supervision: loss hooks for token- or distribution-level supervision, rollout hooks for trajectory metadata, reward hooks for judge or rubric feedback, alignment hooks for token- or span-level mapping, and teacher-sidecar hooks for teacher-side signals such as logits, hidden states, or judge outputs, as used respectively by logit-based, feature-based~\citep{zhang2024dskd}, and rubric-based methods.

This design balances expressiveness and maintainability. Writing every method directly into verl would create method-specific branches, while exposing only a single fixed interface would make alignment-, reward-, or trajectory-based OPD methods difficult to express. Thin hooks keep method logic local while verl remains the shared backend. EasyOPD thus inherits rollout generation, worker orchestration, optimization, and distributed execution from verl; its contribution is the dispatch layer connecting these components to method-local supervision logic.
